\section{增稳控制律设计}
\label{sec:control}

本节基于前文的动力学模型和控制效能分析，设计纵列双发矢量推力构型的增稳与控制律。
控制律采用级联结构，底层为SAS增稳内环，中层为姿态保持/角速度指令外环，
顶层为导航/轨迹规划（不在本文范围内）。

\subsection{级联控制架构}

控制架构分为三个明确的功能层次：

\begin{enumerate}[label=(\arabic*)]
    \item \textbf{操纵输入层}：人类操纵手（通过遥控器）或上层自主控制器
          （轨迹规划器/导航控制器）产生四通道操纵输入——
          油门$\omega_0$（纵向推力）、俯仰指令$\delta_t^{\text{pilot}}$、
          偏航指令$\delta_f^{\text{pilot}}$、滚转指令$\Delta\omega^{\text{pilot}}$；
    \item \textbf{SAS增强层}（Stability Augmentation System）：
          这是内环控制的核心，直接作用于角速率和姿态角反馈。
          SAS的使命是改善短周期阻尼、抑制大气湍流扰动、
          消除配平常值误差（通过积分项），并对飞行员操纵输入进行增稳修正；
    \item \textbf{执行层}：修正后的指令经物理限幅（摆角$\pm\delta_{\max}$，
          差速$\pm\Delta\omega_{\max}$）和电机动态模型后驱动实际执行器。
\end{enumerate}

\textbf{架构特征}：当前设计为\textbf{直接映射反馈增稳}，控制分配采用对角效能矩阵的
直接求逆（式\ref{eq:direct_mapping}），而非伪逆优化或二次规划在线求解。
这一简化依赖于第\ref{sec:allocation}节证明的通道天然解耦性质，
在正交摆轴几何约束下是完全合理的——可视为几何设计简化算法的经典案例。

\begin{figure}[H]
\centering
\begin{tikzpicture}[>=Stealth, node distance=8mm and 12mm,
  box/.style={draw=black!70, fill=black!8, rounded corners=3pt, minimum width=22mm, minimum height=9mm, align=center, font=\footnotesize\sffamily},
  sbox/.style={draw=black!50, fill=black!4, rounded corners=2pt, minimum width=16mm, minimum height=7mm, align=center, font=\footnotesize\sffamily},
  arrow/.style={->, thick, >={Stealth[length=2.5mm]}},
  annot/.style={font=\footnotesize\sffamily},
  dashbox/.style={draw=black!30, densely dashed, rounded corners=6pt, inner sep=6pt},
]
  % ===== 操纵输入（左侧） =====
  \node[box, fill=blue!8, draw=blue!40] (pilot) {操纵输入\\$\omega_0,\delta_t^{\text{pilot}},\delta_f^{\text{pilot}},\Delta\omega^{\text{pilot}}$};
  \node[annot, above=0pt of pilot] {操纵手/自主控制器};

  % ===== SAS层 =====
  \node[box, fill=red!8, draw=red!40, right=of pilot, xshift=4mm] (sas) {SAS 增稳层\\角速率阻尼\\+姿态P+I\\+角速度闭环};

  % ===== 限幅 =====
  \node[sbox, right=of sas, xshift=4mm] (limit) {指令限幅\\$\pm\delta_{\max},\pm\Delta\omega_{\max}$};

  % ===== 执行器 =====
  \node[box, fill=green!8, draw=green!50, right=of limit, xshift=4mm] (act) {执行器\\电机一阶\\滞后+摆座};

  % ===== 推力映射 =====
  \node[box, right=of act, xshift=4mm] (map) {推力矢量\\六维映射\\式(\ref{eq:six_component})};

  % ===== 动力学 =====
  \node[box, fill=orange!8, draw=orange!60, right=of map, xshift=4mm] (dyn) {六自由度\\刚体动力学\\NE+四元数};

  % ===== 传感器反馈（下方） =====
  \node[box, fill=black!5, draw=black!40, below=of sas, yshift=-6mm] (sensor) {传感器反馈\\$\bm{\omega}, \bm{q}, \bm{\theta}$};

  % ===== 扰动 =====
  \node[annot, above=of dyn, yshift=4mm] (dist) {气动扰动/湍流};

  % 前向连接
  \draw[arrow] (pilot) -- (sas);
  \draw[arrow] (sas) -- (limit);
  \draw[arrow] (limit) -- (act);
  \draw[arrow] (act) -- (map);
  \draw[arrow] (map) -- (dyn);
  \draw[arrow] (dist) -- (dyn);

  % 反馈连接
  \draw[arrow] (dyn.south) -- ++(0,-0.4) -| (sensor.east) node[pos=0.6, annot, right=2pt] {状态};
  \draw[arrow] (sensor.west) -| (sas.south) node[pos=0.35, annot, below=2pt] {反馈};

  % SAS 虚线框
  \begin{scope}[on background layer]
    \node[dashbox, fit=(sas)(sensor), label={[annot] above right:{\textbf{SAS 内环}}}] {};
  \end{scope}

  % 输出
  \draw[arrow] (dyn.east) -- ++(0.8,0) node[annot, right] {飞行状态};
\end{tikzpicture}
\caption{纵列双发矢量推力构型级联控制架构。三层结构：操纵输入层生成四通道指令，
SAS增强层通过角速率/姿态反馈修正指令，执行层经物理限幅后驱动电机和摆座。
控制分配采用对角效能矩阵的直接映射（式\ref{eq:direct_mapping}），无需在线优化求解。
内环反馈由IMU提供机体角速率和姿态四元数。}
\label{fig:control_architecture}
\end{figure}


\subsection{SAS基本增稳律}

三个通道的SAS反馈律以效率符号整定为原则。

\textbf{俯仰通道}（效率为负：$\partial M_y / \partial \delta_t = -b T_t < 0$）：
\begin{equation}
    \delta_t^{\text{cmd}} = \delta_t^{\text{pilot}}
    + k_q q + k_\theta \Delta\theta
    + k_i^\theta \int \Delta\theta \, dt
    \label{eq:sas_pitch}
\end{equation}
其中$k_q > 0$（效率为负$\to$正反馈使俯仰速率增大$\to$低头力矩增大$\to$抑制抬头速率）；
$k_\theta > 0$（俯仰角偏离越大$\to$摆角修正越大$\to$恢复力矩越大）；
积分项$k_i^\theta \int \Delta\theta \, dt$消除配平常值误差（如因推力线不完全过质心
或气动配平偏差导致的持续非零俯仰角）。

\textbf{偏航通道}（效率为正：$\partial M_z / \partial \delta_f = a T_f > 0$）：
\begin{equation}
    \delta_f^{\text{cmd}} = \delta_f^{\text{pilot}}
    - k_r r
    \label{eq:sas_yaw}
\end{equation}
效率为正意味着正摆角产生正偏航力矩（右偏航），因此偏航阻尼需要\textbf{负}反馈
（$-k_r r$）——当飞行器向右偏航（$r > 0$），控制律产生负摆角（左偏）
以产生逆方向的偏航力矩来抑制偏航速率。
偏航通道通常不包含姿态积分——航向保持是上层导航闭环的职责，
SAS仅需提供内环阻尼和角速率响应特性的改善。

\textbf{滚转通道}（效率为负：$\partial M_x / \partial \Delta\omega = -2k_Q \omega_0^2 < 0$）：
\begin{equation}
    \Delta\omega^{\text{cmd}} = \Delta\omega^{\text{pilot}}
    + k_p p + k_\phi \Delta\phi
    + k_i^\phi \int \Delta\phi \, dt
    \label{eq:sas_roll}
\end{equation}
效率为负$\to$正反馈，原理同俯仰通道。

\textbf{增益整定指南}（指导性原则，非标定数值）：
\begin{itemize}
    \item \textbf{角速率阻尼增益}（$k_q, k_r, k_p$）：从低值逐步增大，直至角速率阶跃响应
          的非周期阻尼（过阻尼）或目标阻尼比（0.7--1.0）。过高的角速率阻尼增益
          会引入高频噪声放大和执行器饱和风险；
    \item \textbf{姿态比例增益}（$k_\theta, k_\phi$）：通常设定为角速率阻尼增益的
          $\sim 1/4 \cdot \tau_{\text{姿态}}$量级，其中$\tau_{\text{姿态}}$为目标姿态
          响应时间常数；
    \item \textbf{积分增益}（$k_i^\theta, k_i^\phi$）：应足够低，使积分饱和
          恢复时间$\sim 5$--$10$秒。积分增益过高是诱发PIO（驾驶员诱发振荡）的常见原因。
\end{itemize}

\subsection{角速度闭环模式（Rate Command）}

角速度闭环模式（Rate Command / Rate CAC）是一种常见的飞行品质增强模式，
其中操纵杆（或滑块）的偏转直接映射为目标角速度，而非执行器位置指令。
松手后操纵杆弹簧回中、目标角速度归零、飞行器保持当前姿态（速率归零后的姿态）。
这种控制模式在遥控飞行器（如穿越机acro模式）和电传飞行器中广泛应用。

采用纯P控制器（比例控制）追踪目标角速率：
\begin{align}
    \delta_t^{\text{cmd}} &= k_q^{\text{rate}} (q - q_{\text{ref}}),
    &&\partial M_y / \partial \delta_t < 0 \Rightarrow (q - q_{\text{ref}})
    \text{为正时} \delta_t^{\text{cmd}} \text{取正} \\
    \delta_f^{\text{cmd}} &= k_r^{\text{rate}} (r_{\text{ref}} - r),
    &&\partial M_z / \partial \delta_f > 0 \Rightarrow (r - r_{\text{ref}})
    \text{为正时} \delta_f^{\text{cmd}} \text{取负} \\
    \Delta\omega^{\text{cmd}} &= k_p^{\text{rate}} (p - p_{\text{ref}}),
    &&\partial M_x / \partial \Delta\omega < 0 \Rightarrow (p - p_{\text{ref}})
    \text{为正时} \Delta\omega^{\text{cmd}} \text{取正}
    \label{eq:rate_command}
\end{align}

\textbf{性能分析}：纯P控制的角速度闭环存在稳态误差——
为维持非零角速率，需要非零的执行器偏转来克服气动阻尼力矩。
稳态误差可通过积分项消除，但积分项的引入会降低系统相位裕度，
在快速机动中可能导致超调和振荡。对于遥控飞行器（操纵手在闭环中），
适度的稳态误差是可以接受的——操纵手会自然地补偿并通过视觉反馈微调。

角速率闭环的跟踪带宽由P增益和电机响应共同决定。开环传递函数近似为
$G(s) = k^{\text{rate}} \cdot (\partial M / \partial \delta) \cdot
(1/(I s + |L_p|)) \cdot (1/(\tau_m s + 1))$，闭环带宽$\sim k^{\text{rate}}
\cdot |\partial M / \partial \delta| / I$（当电机响应足够快时）。

\subsection{饱和管理与积分器抗饱和}

所有执行器指令须经物理限幅，模拟真实舵机和电机的行程/速率限制：
\begin{equation}
    |\delta_t^{\text{cmd}}| \leq \delta_{\max}, \quad
    |\delta_f^{\text{cmd}}| \leq \delta_{\max}, \quad
    |\Delta\omega^{\text{cmd}}| \leq \Delta\omega_{\max}
    \label{eq:actuator_limits}
\end{equation}

积分器采用独立限幅和条件积分策略：
\begin{itemize}
    \item 积分器输出限幅$|\int \theta| \leq I_\theta^{\max}$防止姿态误差长时间持续
          导致的积分项无限增长（积分饱和，integrator windup）；
    \item 模式切换时积分器归零——当SAS模式从角速率闭环（mode 3）切换回全SAS（mode 1）时，
          旧模式下累积的积分值（基于不同误差信号）直接清零，防止在新模式刚启动时产生
          瞬态冲击（此问题在工程实践中被发现并修复）。
\end{itemize}

\subsection{差速分配与最大转速保护}

$\Delta\omega^{\text{cmd}}$由式(\ref{eq:diff_def})的平方根分配产生两电机目标转速。
此分配具有保持$\omega_f^2 + \omega_t^2 = 2\omega_0^2$常数、使一阶近似总推力不变的
优越性质。

然而在极限工况（高油门+大差速）下，一侧电机的目标转速可能超过最大允许转速$\omega_{\max}$。
在此情况下，必须施加电机级硬限幅：
\begin{equation}
    \omega_f^{\text{target}} \leftarrow \min(\omega_f^{\text{target}}, \omega_{\max}), \quad
    \omega_t^{\text{target}} \leftarrow \min(\omega_t^{\text{target}}, \omega_{\max})
    \label{eq:max_clamp}
\end{equation}

转速硬限幅的副作用是：\textbf{被钳位的电机转速低于目标值，导致平方转速和偏离$2\omega_0^2$，
总推力低于预期}。此外，由于只有一侧电机被钳位，差速指令的滚转力矩也不再精确等于
$-2k_Q\omega_0^2\Delta\omega^{\text{cmd}}$。SAS应在遥测中标记$w_{\max}$饱和状态，
提醒操纵手执行器已达物理极限，当前的控制指令无法被完全执行。


\section{讨论：构型优势、局限与扩展路径}
